\begin{thanks}
  论文写到这里的时候，学科楼前的海棠花快要落完了。风一吹，花瓣像轻轻的标点，落在台阶上、走廊里，也落在这一段时光的末尾。春天走到尾声，我的求学生涯也慢慢收拢成一个句号。若要用一句话概括此刻的心情，我最想说的不是“完成了”，而是“谢谢”。

  首先要感谢我的导师张伯雷老师。谢谢您在我最迷茫、最容易怀疑自己的时候，仍愿意把时间分给我，把问题重新摆到光亮处：该从哪里下手、为什么这样做、怎样把一件事做得更扎实。您对研究的严谨与克制、对细节的敏感，以及在关键节点上不动声色的托举，构成了我这段经历里最可靠的坐标。

  感谢课题组的同学们。许多想法并不是在某个瞬间“突然拥有”，而是在一次次讨论、争辩与推演里逐渐成形：被追问到语塞的那几分钟、在白板前反复擦写的那些式子、把一组不理想的结果拆开重看的耐心。谢谢你们在我卡住的时候一起把思路捋直，在我实验崩掉的时候陪我把错误找出来，也在我有一点点进展的时候，愿意认真听我把它讲完。

  感谢陪我一起吃饭散步的朋友。你们未必关心我用的是哪一种算法，却会关心我有没有好好吃饭、能不能按时睡觉；你们不一定看得懂图表和符号，却愿意在我最焦虑的时候陪我把情绪放回原位。那些看似与论文无关的温柔与照拂，恰恰让我有力气继续写下去。

  感谢美丽的南京邮电大学。谢谢这几年里一路相伴的花花草草：雨后更亮的叶子、傍晚被夕阳照得发暖的树影、以及每一个在图书馆和实验室之间奔波的清晨与深夜。也谢谢校园里那些可爱的小猫——它们不懂我的焦虑与截止日期，却用一次擦肩、一次停留，让我在最紧绷的时候重新想起呼吸和松弛。很多时候，生活里这些微小而确切的美，恰好把我从无数次“差一点坚持不下去”里拉回来。

  最后，也想谢谢那个在反复修改、不断推翻重来的过程中没有放弃的自己。愿这段经历留下的，不止是一份阶段性的成果，更是一种在迷雾里继续前行的勇气与能力。

\end{thanks}
